\section{Ambient tripartite graph versus actual incidence}
\label{sec:ambient-graph}

Let \(K\) index the nonprimitive kernel sites. Write
\[
 B:\C^J\to\C^K,\qquad
 A:\C^K\to\C^I,\qquad
 C=AB.
\]
For real-linear maps, \(A_{ik}B_{kj}\) below means composition of the
corresponding real coordinate blocks, and the sum is taken in
\(\operatorname{Hom}_{\mathbb R}(\C,\C)\). In the pair-dependent
TPC--76 specialization all maps are complex-linear and these are the
usual complex scalar entries.

\begin{definition}[Ambient tripartite graph]
\label{def:tripartite}
The graph \(\cG_{A,B}\) has vertex set
\[
 J\sqcup K\sqcup I
\]
and edges
\[
 j\sim k\iff B_{kj}\ne0,
 \qquad
 k\sim i\iff A_{ik}\ne0.
\]
\end{definition}

\begin{theorem}[Safe coarse factorization and refinement]
\label{thm:ambient-refinement}
If \(i\in I\) and \(j\in J\) lie in different connected components of
\(\cG_{A,B}\), then \(C_{ij}=0\). Consequently, the partition induced
on \(I\sqcup J\) by \(\cG_{A,B}\) gives a valid block decomposition of
\(C\). The actual graph \(\cG_C\) refines this partition and can split
an ambient component further.
\end{theorem}

\begin{proof}
Each summand in
\[
 C_{ij}
 =
 \sum_{k\in K}A_{ik}B_{kj}
\]
can be nonzero only when there is a length-two path
\(j-k-i\) in \(\cG_{A,B}\). If \(i,j\) are in different components,
every summand vanishes. Thus no actual edge crosses the ambient
partition. The converse can fail because the nonzero summands can
cancel, so \(\cG_C\) may contain fewer edges.
\end{proof}

\begin{example}[Exact cancellation splits an ambient block]
\label{ex:ambient-cancellation}
Take two sites, one analysis row, and one filter variable, with
\[
 B=
 \begin{pmatrix}1\\1\end{pmatrix},
 \qquad
 A=
 \begin{pmatrix}1&-1\end{pmatrix}.
\]
The tripartite graph is connected, but \(C=AB=0\). The column is
cost-null and the row is immutable. A graph built before multiplication
misses this exact deletion.
\end{example}

\begin{remark}[Unknown is not zero]
When \(C_{ij}\) is known only inside an interval containing both zero
and nonzero values, the edge must be retained. Deleting it by a
tolerance can manufacture a false component factorization and a false
dual certificate.
\end{remark}

\subsection{Assembly order}

For the TPC application, the safe order is:
\begin{enumerate}[label=\textup{(\arabic*)},leftmargin=2.8em]
 \item freeze the literal \(q\)-resolved carrier and raw normalization;
 \item assemble \(B_{\boldsymbol L}\) while retaining every declared
       \((d,q)\) provenance column; exact complete-key cancellation may
       occur inside one fixed \(q\)-slice, but distinct \(q\)-columns
       are not merged;
 \item assemble the full low-plaquette and exact-tail operator \(A_D\);
 \item compute or validate \(C=A_DB_{\boldsymbol L}\);
 \item delete exact zeros and build \(\cG_C\); and
 \item solve and certify each resulting component.
\end{enumerate}
Changing this order can confuse ambient support, actual support, and
post-\(q\) aggregation.
